\documentclass[11pt,a4paper]{article}
\usepackage[italian]{babel}
\usepackage[margin=2.5cm]{geometry}
\usepackage{parskip}
\usepackage{amsmath, amssymb, amsthm}
\usepackage{booktabs}
\usepackage{array}
\usepackage{xcolor}
\usepackage{needspace}
\usepackage{enumitem}
\usepackage{listings}
\usepackage{tikz}
\usepackage[most]{tcolorbox}
\usepackage[hidelinks]{hyperref}
\usetikzlibrary{decorations.pathreplacing, shapes.geometric, arrows.meta, positioning}

% Niente vedove/orfane: i paragrafi non lasciano righe isolate a fine/inizio pagina
\widowpenalty=10000
\clubpenalty=10000

% Elenchi più compatti
\setlist{itemsep=2pt, parsep=0pt, topsep=4pt, partopsep=0pt}

% --- Riquadri con bordo nero, senza numero (si spezzano tra due pagine) ---
% Uso: \begin{definizione} ... \end{definizione}
%      \begin{teorema}[Nome] ... \end{teorema}  (il nome è facoltativo)
\tcbset{nota/.style={enhanced jigsaw, breakable, colback=white,
  colframe=black, boxrule=0.5pt, sharp corners}}
\NewTColorBox{definizione}{}{nota,
  before upper={\textbf{Definizione.}\ }}
\NewTColorBox{teorema}{o}{nota,
  before upper={\textbf{Teorema\IfValueT{#1}{ (#1)}.}\ }}
\NewTColorBox{proposizione}{o}{nota,
  before upper={\textbf{Proposizione\IfValueT{#1}{ (#1)}.}\ }}
\NewTColorBox{proprieta}{o}{nota,
  before upper={\textbf{Proprietà\IfValueT{#1}{ (#1)}.}\ }}
\NewTColorBox{assioma}{o}{nota, boxrule=1pt,
  before upper={\textbf{Assioma\IfValueT{#1}{ (#1)}.}\ }}

% --- Ambienti semplici (senza riquadro) ---
% Il titolo sta in cima al blocco, anche se il contenuto inizia con un riquadro o
% con codice; e non resta mai da solo in fondo a una pagina.
% Uso: \begin{esempio}[Titolo facoltativo] ... \end{esempio}
\newcommand{\newrunin}[2]{%
  \NewDocumentEnvironment{#1}{o}{%
    \par\Needspace{4\baselineskip}\addvspace{\medskipamount}\noindent
    \textbf{#2}\IfValueT{##1}{ (##1)}\textbf{.}\ \ignorespaces}%
  {\par\addvspace{\medskipamount}}}
\newrunin{esempio}{Esempio}
\newrunin{esercizio}{Esercizio}
\newrunin{osservazione}{Osservazione}

% V e F nelle tavole di verità
\newcommand{\V}{\textsf{V}}
\newcommand{\F}{\textsf{F}}

% --- Codice C ---
% Uso: \begin{codice} ... \end{codice}      codice C numerato
%      \begin{comando} ... \end{comando}    comandi da terminale
%      \begin{risultato} ... \end{risultato} output di un programma
%      \lstinline|printf("ciao");|           codice dentro al testo
\lstdefinestyle{c}{
  language=C,
  basicstyle=\ttfamily\small,
  keywordstyle=\color{blue}\bfseries,
  commentstyle=\color{gray}\itshape,
  stringstyle=\color{red!70!black},
  numbers=left,
  numberstyle=\tiny\color{gray},
  numbersep=8pt,
  columns=flexible,
  keepspaces=true,
  breaklines=true,
  showstringspaces=false,
  tabsize=4,
  morekeywords={include, define},
  % lettere accentate nei commenti
  literate={à}{{\`a}}1 {è}{{\`e}}1 {é}{{\'e}}1 {ì}{{\`i}}1
           {ò}{{\`o}}1 {ù}{{\`u}}1 {È}{{\`E}}1 {À}{{\`A}}1
}
\lstset{style=c}
\newtcblisting{codice}{listing only, nota, left=6mm,
  listing options={style=c}}
\newtcblisting{comando}{listing only, nota,
  listing options={basicstyle=\ttfamily\small, numbers=none, language={},
    columns=flexible, keepspaces=true, breaklines=true}}
\newtcblisting{risultato}{listing only, enhanced jigsaw, breakable,
  colback=black!5, colframe=black, boxrule=0.5pt, sharp corners,
  listing options={basicstyle=\ttfamily\small, numbers=none, language={},
    columns=flexible, keepspaces=true, breaklines=true}}

% --- Diagrammi di flusso (simboli standard) ---
\tikzset{
  terminale/.style = {ellipse, draw, minimum width=2.5cm, minimum height=0.9cm},
  io/.style        = {trapezium, trapezium left angle=70, trapezium right angle=110,
                      draw, minimum width=2.5cm, minimum height=0.9cm},
  processo/.style  = {rectangle, draw, minimum width=2.5cm, minimum height=0.9cm},
  decisione/.style = {diamond, draw, aspect=2, minimum width=2.5cm},
  freccia/.style   = {thick, -{Stealth}}
}

\title{Matematica discreta\\[0.4em]\large Appunti}
\author{Marco Cerbella}
\date{Università degli Studi di Perugia -- A.A. 2026/2027\\ Aggiornato al \today}

\begin{document}
\maketitle
\setcounter{tocdepth}{1}
\tableofcontents
\clearpage

\section[Logica (Lezione 1 -- 21 settembre 2026)]{Logica}
\noindent\emph{Lezione 1 -- 21 settembre 2026}
\subsection{Enunciati e connettivi}

\begin{definizione}
Un \emph{enunciato} (o \emph{proposizione}) è una frase di cui possiamo
dire con certezza se è \V\ (vera) oppure \F\ (falsa). Gli enunciati si
indicano con $p, q, r, \dots$
\end{definizione}

\begin{esempio}
Sono enunciati: ``$5 > 3$'', ``$99$ è divisibile per $3$'',
``$508$ è multiplo di $5$''. Invece ``$n$ è un numero primo'' \textbf{non}
è un enunciato, perché la sua verità dipende da $n$.
\end{esempio}

Dati due enunciati $p, q$ se ne costruiscono di nuovi con i
\emph{connettivi}:
\begin{center}
\begin{tabular}{cll}
$p \land q$ & ``$p$ e $q$'' & congiunzione \\
$p \lor q$ & ``$p$ o $q$'' & disgiunzione \\
$\neg p$ & ``non $p$'' & negazione \\
$p \to q$ & ``se $p$ allora $q$'' & implicazione \\
$p \leftrightarrow q$ & ``$p$ se e soltanto se $q$'' & coimplicazione
\end{tabular}
\end{center}

\subsubsection*{Tavola di verità}

\begin{center}
\begin{tabular}{cc|c|c|c|c}
$p$ & $q$ & $p \land q$ & $p \lor q$ & $p \to q$ & $p \leftrightarrow q$ \\ \hline
\V & \V & \V & \V & \V & \V \\
\V & \F & \F & \V & \F & \F \\
\F & \V & \F & \V & \V & \F \\
\F & \F & \F & \F & \V & \V
\end{tabular}
\end{center}

\subsection{Proprietà dei connettivi}

\begin{proposizione}
Siano $p, q, r$ enunciati. Allora:
\begin{enumerate}
    \item $p \lor q \iff q \lor p$;
    \item $p \land q \iff q \land p$;
    \item $(p \land q) \lor r \iff (p \lor r) \land (q \lor r)$;
    \item $(p \lor q) \land r \iff (p \land r) \lor (q \land r)$;
    \item $\neg(p \land q) \iff (\neg p) \lor (\neg q)$;
    \item $\neg(p \lor q) \iff (\neg p) \land (\neg q)$;
    \item $p \to q \iff (\neg p) \lor q$;
    \item $p \to q \iff \neg q \to \neg p$;
    \item $p \leftrightarrow q \iff (p \to q) \land (q \to p)$.
\end{enumerate}
\end{proposizione}

Due enunciati composti sono equivalenti ($\iff$) quando hanno la stessa
tavola di verità. Le prime due proprietà sono la commutatività di $\lor$ e
$\land$; la 3 e la 4 sono le proprietà distributive; la 5 e la 6 sono le
leggi di De Morgan; la 8 è la \emph{contronominale}.

\begin{proof}[Verifica della proprietà 7]
Confrontiamo le tavole di verità di $p \to q$ e di $(\neg p) \lor q$:
\begin{center}
\begin{tabular}{cc|c|c|c}
$p$ & $q$ & $p \to q$ & $\neg p$ & $(\neg p) \lor q$ \\ \hline
\V & \V & \V & \F & \V \\
\V & \F & \F & \F & \F \\
\F & \V & \V & \V & \V \\
\F & \F & \V & \V & \V
\end{tabular}
\end{center}
Le colonne di $p \to q$ e di $(\neg p) \lor q$ sono uguali.
\end{proof}

\begin{proof}[Verifica della proprietà 4]
Servono otto righe, una per ogni combinazione dei valori di $p, q, r$:
\begin{center}
\begin{tabular}{ccc|c|c|c|c|c}
$p$ & $q$ & $r$ & $p \lor q$ & $(p \lor q) \land r$ & $p \land r$ &
$q \land r$ & $(p \land r) \lor (q \land r)$ \\ \hline
\V & \V & \V & \V & \V & \V & \V & \V \\
\V & \V & \F & \V & \F & \F & \F & \F \\
\V & \F & \V & \V & \V & \V & \F & \V \\
\V & \F & \F & \V & \F & \F & \F & \F \\
\F & \V & \V & \V & \V & \F & \V & \V \\
\F & \V & \F & \V & \F & \F & \F & \F \\
\F & \F & \V & \F & \F & \F & \F & \F \\
\F & \F & \F & \F & \F & \F & \F & \F
\end{tabular}
\end{center}
Le colonne di $(p \lor q) \land r$ e di $(p \land r) \lor (q \land r)$
coincidono. Le altre proprietà si verificano allo stesso modo.
\end{proof}

\subsection{Formule aperte e quantificatori}

\begin{definizione}
Una \emph{formula aperta} è un'affermazione che contiene un'incognita, un
parametro o una variabile.
\end{definizione}

\begin{esempio}
\[
p(x):\ x + 3 = 0, \qquad
q(n):\ \text{``$n$ è un numero pari''}, \qquad
r(k):\ \text{``$k$ è multiplo di $5$''}.
\]
\end{esempio}

Una formula aperta non è né vera né falsa. Per passare dalla formula aperta
a un enunciato ci sono tre modi.

\begin{enumerate}
    \item \textbf{Sostituzione}: si sostituisce alla variabile un valore
          preciso. Ad esempio $p(3)$ è l'enunciato $3 + 3 = 0$ (falso),
          $p(-3)$ è $-3 + 3 = 0$ (vero), $q(8)$ è ``$8$ è un numero pari''
          (vero).

    \item \textbf{Quantificatore esistenziale} $\exists$ (``esiste''):
          \[
          \exists\, x \in \mathbb{R} : p(x)
          \]
          si legge ``esiste $x$ reale tale che $x + 3 = 0$'' ed è vero
          (basta $x = -3$). Invece
          $\exists\, x \in \mathbb{R} : x^2 = -10$ è falso.

    \item \textbf{Quantificatore universale} $\forall$ (``per ogni''):
          $\forall\, x \in \mathbb{R} : x + 3 = 0$ è falso, mentre
          $\forall\, x \in \mathbb{R} : x^2 \geq 0$ è vero.
\end{enumerate}

\begin{proposizione}[Negazione dei quantificatori]
\[
\neg\big(\forall x : p(x)\big) \iff \exists\, x : \neg p(x), \qquad
\neg\big(\exists\, x : p(x)\big) \iff \forall x : \neg p(x).
\]
Il secondo enunciato si scrive anche $\nexists\, x : p(x)$.
\end{proposizione}

Per negare una frase con un quantificatore si scambia $\forall$ con
$\exists$ e si nega la formula.

\subsection{Implicazione: condizioni necessarie e sufficienti}

L'implicazione $p \to q$ si legge in tre modi:
\begin{itemize}
    \item ``se $p$ allora $q$'';
    \item ``$p$ è condizione \emph{sufficiente} per $q$'';
    \item ``$q$ è condizione \emph{necessaria} per $p$''.
\end{itemize}

\begin{esempio}
Per ogni $n \in \mathbb{N}$: se $n$ è divisibile per $4$ (è $p$), allora
$n$ è pari (è $q$). Essere divisibile per $4$ è sufficiente per essere
pari; essere pari è necessario per essere divisibile per $4$.
\end{esempio}

\section[Insiemi (Lezione 1 -- 21 settembre 2026)]{Insiemi}
\noindent\emph{Lezione 1 -- 21 settembre 2026}
\subsection{Insiemi ed elementi}

\begin{definizione}
Un \emph{insieme} è una collezione (o famiglia) di oggetti di cui possiamo
dire con precisione se un oggetto appartiene o no alla famiglia. Un oggetto
di un insieme si chiama \emph{elemento}.
\end{definizione}

Gli insiemi si indicano di solito con lettere maiuscole ($A, B, X, S, V,
\dots$), gli elementi con lettere minuscole ($x, y, v, k, \dots$).
\begin{itemize}
    \item $x \in S$: l'elemento $x$ appartiene all'insieme $S$;
    \item $y \notin S$: $y$ non appartiene a $S$.
\end{itemize}

Un insieme si può dare in due modi.
\begin{enumerate}
    \item \textbf{Per elencazione}: si scrivono gli elementi tra graffe.
          Ad esempio $V = \{1, 3, 5, 88\}$ definisce l'insieme $V$ per
          elencazione.
    \item \textbf{Per descrizione} (o per caratteristica): ad esempio
          $X = \{\text{le lettere dell'alfabeto inglese}\}$.
\end{enumerate}

Quando la descrizione usa una formula aperta si scrive
\[
W = \{x \mid x^2 - 3x + 2 = 0\}, \qquad
A = \{x \in \mathbb{R} \mid P(x)\}, \qquad
B = \{x \in \mathbb{R} \mid x > 3\} = \{x : x > 3\},
\]
dove ``$\mid$'' (oppure ``$:$'') si legge ``tale che''.

Un insieme con un solo elemento, come $U = \{4\}$ (per cui $4 \in U$ e
$3 \notin U$), si chiama \emph{singoletto} (in inglese \emph{singleton}).

L'\emph{insieme vuoto}, indicato con $\emptyset = \{\ \}$, è l'insieme senza
elementi: per ogni $x$ vale $x \notin \emptyset$ (cioè
$\nexists\, x : x \in \emptyset$).

\subsection{Sottoinsiemi e uguaglianza}

\begin{definizione}
Siano $A, B$ insiemi. $B$ è un \emph{sottoinsieme} di $A$ se
\[
\forall\, x \in B \implies x \in A.
\]
Si scrive $B \subset A$ oppure $B \subseteq A$; equivalentemente
$A \supseteq B$ (o $A \supset B$).
\end{definizione}

\begin{definizione}
Due insiemi $S, T$ sono \emph{uguali}, $S = T$, se hanno gli stessi
elementi, cioè
\[
x \in S \iff x \in T.
\]
\end{definizione}

\begin{osservazione}
$S = T$ se e solo se $S \subseteq T$ e $T \subseteq S$.
\end{osservazione}

Dato un insieme $A$ valgono sempre $\emptyset \subseteq A$ e
$A \subseteq A$.

\subsubsection*{Insiemi numerici}

\begin{center}
\begin{tabular}{ll}
$\mathbb{N} = \{0, 1, 2, 3, \dots\}$ & numeri naturali \\
$\mathbb{Z} = \{0, \pm 1, \pm 2, \pm 3, \dots\}
              = \{\dots, -3, -2, -1, 0, 1, 2, \dots\}$ & numeri interi \\
$\mathbb{Q} = \left\{ \dfrac{a}{b} \;\middle|\; a, b \in \mathbb{Z},\
              b \neq 0 \right\}$ & numeri razionali \\
$\mathbb{R}$ & numeri reali \\
$\mathbb{C}$ & numeri complessi
\end{tabular}
\end{center}

Valgono le inclusioni
\[
\mathbb{N} \subseteq \mathbb{Z} \subseteq \mathbb{Q} \subseteq
\mathbb{R} \subseteq \mathbb{C}.
\]
Con $i$ si indica l'unità immaginaria, $i^2 = -1$. Ad esempio:
\begin{center}
\begin{tabular}{llll}
$-2 \in \mathbb{Z}$, & $\frac12 \in \mathbb{Q}$, &
$\sqrt{2}, \pi \in \mathbb{R}$, & $i \in \mathbb{C}$, \\
$-2 \notin \mathbb{N}$, & $\frac12 \notin \mathbb{Z}$, &
$\sqrt{2}, \pi \notin \mathbb{Q}$, & $i \notin \mathbb{R}$.
\end{tabular}
\end{center}

\subsection{Insieme delle parti}

\begin{definizione}
Sia $A$ un insieme. L'\emph{insieme delle parti} di $A$ è l'insieme di
tutti i sottoinsiemi di $A$:
\[
\mathcal{P}(A) = \{B \mid B \subseteq A\}.
\]
\end{definizione}

\begin{osservazione}
Per ogni insieme $A$: $\emptyset \subseteq A$ e $\emptyset \in
\mathcal{P}(A)$; $A \subseteq A$ e $A \in \mathcal{P}(A)$.
\end{osservazione}

\begin{esempio}
Sia $A = \{1, 2, 3\}$. Allora
\[
\mathcal{P}(A) = \big\{\emptyset,\ \{1\},\ \{2\},\ \{3\},\
\{1,3\},\ \{2,3\},\ \{1,2\},\ A\big\}.
\]
Bisogna distinguere bene appartenenza e inclusione:
\begin{itemize}
    \item $1 \in A$, \quad $\{1\} \subseteq A$, \quad
          $\{1, 3\} \subseteq A$;
    \item $\{1\} \in A$: \textbf{no}; \quad $\{1\} \in \mathcal{P}(A)$:
          sì;
    \item $A \subseteq \mathcal{P}(A)$: \textbf{no}; \quad
          $A \in \mathcal{P}(A)$: sì.
\end{itemize}
\end{esempio}

\begin{esempio}
$\mathcal{P}(\emptyset) = \{\emptyset\}$ e
\[
\mathcal{P}(\{\emptyset\}) = \mathcal{P}(\mathcal{P}(\emptyset))
= \{\emptyset, \{\emptyset\}\}.
\]
\end{esempio}

\subsection{Operazioni tra insiemi}

Siano $A, B$ insiemi.

\begin{definizione}
\begin{itemize}
    \item \textbf{Unione}: $A \cup B = \{x \mid x \in A \lor x \in B\}$.
    \item \textbf{Intersezione}: $A \cap B = \{x \mid x \in A \land
          x \in B\}$.
    \item \textbf{Differenza} tra $A$ e $B$:
          $A \setminus B = \{x \mid x \in A \text{ e } x \notin B\}$.
    \item \textbf{Differenza simmetrica} tra $A$ e $B$:
          $A \,\triangle\, B = (A \setminus B) \cup (B \setminus A)$.
\end{itemize}
$A$ e $B$ sono \emph{disgiunti} se $A \cap B = \emptyset$.
\end{definizione}

Esercizio: $A \,\triangle\, B = (A \cup B) \setminus (A \cap B)$.

\begin{definizione}
Se $A \subseteq U$, il \emph{complementare} di $A$ in $U$ è l'insieme
\[
A^c = U \setminus A.
\]
\end{definizione}

\begin{esempio}
Siano $A = \{1, 2, 3, 4\} = \{4, 1, 3, 2\}$ (l'ordine non conta) e
$B = \{1, 3, 6, 9, 12\}$. Allora
\begin{align*}
A \cup B &= \{1, 2, 3, 4, 6, 9, 12\}, & A \cap B &= \{1, 3\}, \\
A \setminus B &= \{2, 4\}, & B \setminus A &= \{6, 9, 12\}, \\
A \,\triangle\, B &= \{2, 4, 6, 9, 12\}.
\end{align*}
\end{esempio}

\subsection{Prodotto cartesiano}

\begin{definizione}
Siano $A, B$ insiemi. Il \emph{prodotto cartesiano} di $A$ e $B$ è
\[
A \times B = \{(a, b) \mid a \in A,\ b \in B\},
\]
l'insieme delle \emph{coppie ordinate}. Vale $(a, b) \neq (b, a)$ in
generale, e
\[
(x, y) = (x', y') \iff x = x' \text{ e } y = y'.
\]
\end{definizione}

\begin{esempio}
Siano $A = \{1, 2, 3\}$ e $B = \{v, w\}$. Allora
\begin{align*}
A \times B &= \{(1,v), (1,w), (2,v), (2,w), (3,v), (3,w)\}, \\
B \times A &= \{(v,1), (w,1), (v,2), (w,2), (v,3), (w,3)\},
\end{align*}
quindi $A \times B \neq B \times A$. Inoltre
\[
A \times A = A^2 = \{(1,1), (1,2), (1,3), (2,1), (2,2), (2,3), (3,1),
(3,2), (3,3)\}.
\]
\end{esempio}

Analogamente
\[
A \times B \times C = \{(a, b, c) \mid a \in A,\ b \in B,\ c \in C\},
\qquad A \times A \times A = A^3.
\]

\section[Relazioni di equivalenza (Lezione 2 -- 22 settembre 2026)]{Relazioni di equivalenza}
\noindent\emph{Lezione 2 -- 22 settembre 2026}
\subsection{Relazioni}

\begin{definizione}
Una \emph{relazione} $R$ (o $\rho$) su un insieme $A$ è un sottoinsieme del
prodotto cartesiano $A \times A$:
\[
R \subseteq A \times A.
\]
Se $(a, b) \in R$ si scrive $a \mathrel{R} b$.
\end{definizione}

\begin{esempio}
Sia $A = \{a, b, c\}$ e $R = \{(a,a), (a,b), (c,b)\}$. Si ha
$a \mathrel{R} b$, mentre $(b, a) \notin R$, cioè $b \not\mathrel{R} a$.
Questa relazione è transitiva, non è simmetrica e non è riflessiva
($b \not\mathrel{R} b$).
\end{esempio}

\begin{esempio}[Relazioni su $\mathbb{N}$ e $\mathbb{R}$]
\leavevmode
\begin{enumerate}
    \item Su $\mathbb{N}$ la relazione $\leq$: $(5, 8) \in\ \leq$ perché
          $5 \leq 8$, mentre $(10, 3) \notin\ \leq$ perché $10 \nleq 3$.
          È riflessiva ($a \leq a$) e transitiva
          ($a \leq b \land b \leq c \Rightarrow a \leq c$). Non è
          simmetrica: $3 \leq 5$ ma $5 \nleq 3$.

    \item Su $\mathbb{R}$ sia
          $R_1 = \{(x, y) \in \mathbb{R}^2 \mid x + 2y = 3\}
          \subseteq \mathbb{R} \times \mathbb{R}$.
          Si ha $(2, 2) \notin R_1$ perché $2 + 2 \cdot 2 = 6 \neq 3$;
          invece $1 \mathrel{R_1} 1$ (perché $1 + 2 = 3$), ma
          $2 \not\mathrel{R_1} 2$: $R_1$ \textbf{non è riflessiva}.
          Inoltre $(-1, 2) \in R_1$ perché $-1 + 2 \cdot 2 = 3$, ma
          $(2, -1) \notin R_1$ perché $2 + 2 \cdot (-1) = 0 \neq 3$:
          $R_1$ non è simmetrica.

    \item Su $\mathbb{N}$: $h \mathrel{R_2} k \iff h - k = 12$. Allora
          \[
          R_2 = \{(12, 0), (13, 1), (14, 2), \dots\}
              = \{(12 + k, k) \mid k \in \mathbb{N}\}.
          \]
          $(0, 0) \notin R_2$: non è riflessiva.

    \item Su $\mathbb{N}$: $m \mathrel{R_3} n \iff m + n = 5$. Allora
          \[
          R_3 = \{(0,5), (1,4), (2,3), (3,2), (4,1), (5,0)\}.
          \]
          Non è riflessiva ($0 \not\mathrel{R_3} 0$). È simmetrica.
          \textbf{Non} è transitiva: $1 \mathrel{R_3} 4$ e
          $4 \mathrel{R_3} 1$, ma $1 \not\mathrel{R_3} 1$.
\end{enumerate}
\end{esempio}

\begin{definizione}
Se $R$ è una relazione su $A$, la \emph{relazione opposta} $R^{op}$ si
ottiene scambiando di posto gli elementi delle coppie di $R$:
\[
(a, b) \in R \iff (b, a) \in R^{op}, \qquad
a \mathrel{R} b \iff b \mathrel{R^{op}} a.
\]
\end{definizione}

\begin{definizione}
Sia $R$ una relazione su $A$. Si dice che $R$ è:
\begin{itemize}
    \item \emph{riflessiva} se $\forall\, a \in A \quad a \mathrel{R} a$;
    \item \emph{simmetrica} se $a \mathrel{R} b \Rightarrow
          b \mathrel{R} a$;
    \item \emph{transitiva} se $a \mathrel{R} b \land b \mathrel{R} c
          \Rightarrow a \mathrel{R} c$.
\end{itemize}
\end{definizione}

\subsection{Relazioni di equivalenza}

\begin{definizione}
Una relazione $R$ su $A$ riflessiva, simmetrica e transitiva si dice
\emph{relazione di equivalenza}.
\end{definizione}

\begin{esempio}
\leavevmode
\begin{enumerate}
    \item Dato un insieme $A$, la relazione $R = A \times A$ è di
          equivalenza (tutti sono in relazione con tutti). Per ogni
          $a \in A$ si ha $[a] = A$, quindi $A/_R = \{[a]\} = \{A\}$.

    \item In $\mathbb{N}$ la relazione di uguaglianza,
          $n \mathrel{R} k \iff n = k$, è di equivalenza. Le classi sono
          $[n]_= = \{n\}$, ad esempio $[5]_= = \{5\}$, e
          $\mathbb{N}/_= = \{[0], [1], [2], \dots\}$.

    \item Sia $S$ l'insieme delle rette di un piano. Due rette
          $r, s \in S$ sono in relazione se $r$ è parallela a $s$,
          $r \parallel s$. È riflessiva ($s \parallel s$), simmetrica
          ($r \parallel s \Rightarrow s \parallel r$) e transitiva
          ($r \parallel s \land s \parallel t \Rightarrow r \parallel t$):
          è una relazione di equivalenza. La classe di $r$ è
          $[r]_{\parallel} = \{s \mid s \parallel r\}$, e se
          $r_1 \parallel r_2$ allora $[r_1]_{\parallel} =
          [r_2]_{\parallel}$.

    \item Sia $A = \{1, 2, 3, 4, 5, 6\}$ e
          \[
          \begin{split}
          R = \{&(1,1), (1,2), (1,6), (2,1), (2,6), (2,2), (6,6),\\
          &(6,1), (6,2), (3,3), (4,4), (4,5), (5,4), (5,5)\}.
          \end{split}
          \]
          È riflessiva e simmetrica. È transitiva: ad esempio
          $(1,6), (6,2) \in R \Rightarrow (1,2) \in R$ e
          $4 \mathrel{R} 5, 5 \mathrel{R} 4 \Rightarrow 4 \mathrel{R} 4$.
          Quindi è di equivalenza, con
          \[
          [1]_R = \{1, 2, 6\} = [2]_R = [6]_R, \quad [3]_R = \{3\},
          \quad [4]_R = [5]_R = \{4, 5\},
          \]
          e $A/_R = \{[1]_R, [3]_R, [4]_R\}$.
\end{enumerate}
\end{esempio}

\begin{esempio}[Relazioni che non sono di equivalenza]
\leavevmode
\begin{enumerate}
    \setcounter{enumi}{4}
    \item Su $\mathbb{Z}$: $k \mathrel{R} h \iff h + k = 10$. Non è
          riflessiva ($0 \not\mathrel{R} 0$), quindi non è di equivalenza.

    \item $A = \{a, b, c\}$ e $R = \{(a,a), (b,b), (a,c), (c,a)\}$.
          Non è riflessiva ($c \not\mathrel{R} c$). È simmetrica. Non è
          transitiva: $c \mathrel{R} a$ e $a \mathrel{R} c$, ma
          $c \not\mathrel{R} c$. (Altri casi, come
          $c \mathrel{R} a, a \mathrel{R} a \Rightarrow c \mathrel{R} a$,
          vanno bene: basta un caso che fallisce.)

    \item Su $\mathcal{P}(A)$ la relazione $X \subseteq Y$ (``essere
          sottoinsieme di'', con $X, Y \in \mathcal{P}(A)$) non è
          simmetrica.

    \item Su $\mathbb{N}$ la relazione $\leq$ non è di equivalenza: non è
          simmetrica.
\end{enumerate}
\end{esempio}

\subsection{Classi di equivalenza e insieme quotiente}

\begin{definizione}
Sia $A$ un insieme e $\rho$ una relazione di equivalenza su $A$. La
\emph{classe di equivalenza} di $a \in A$ è l'insieme
\[
[a]_\rho = \{x \in A \mid x \mathrel{\rho} a\}.
\]
L'insieme di tutte le classi di equivalenza è detto \emph{insieme
quotiente} di $A$ rispetto a $\rho$:
\[
A/_\rho = \{[a]_\rho \mid a \in A\}.
\]
\end{definizione}

\begin{osservazione}
Sia $\rho$ una relazione di equivalenza su $A$ e $a \in A$.
\begin{enumerate}
    \item $[a]_\rho \subseteq A$, quindi $[a]_\rho \in \mathcal{P}(A)$ e
          $A/_\rho = \{[a]_\rho \mid a \in A\} \subseteq \mathcal{P}(A)$.
    \item $[a]_\rho = [b]_\rho \iff a \mathrel{\rho} b \iff
          [a]_\rho \cap [b]_\rho \neq \emptyset$.
    \item Se $\rho$ è di equivalenza, $\rho^{op} = \rho$.
\end{enumerate}
\end{osservazione}

\subsection{Partizioni}

\begin{definizione}
Sia $A$ un insieme e $\mathcal{F} \subseteq \mathcal{P}(A)$ una famiglia di
sottoinsiemi di $A$. $\mathcal{F}$ è una \emph{partizione} di $A$ se:
\begin{enumerate}
    \item $\forall\, X \in \mathcal{F} \quad X \neq \emptyset$;
    \item $\forall\, X, Y \in \mathcal{F} \quad X \neq Y \Rightarrow
          X \cap Y = \emptyset$;
    \item $\displaystyle \bigcup_{X \in \mathcal{F}} X = A$.
\end{enumerate}
\end{definizione}

In parole: i pezzi non sono vuoti, sono a due a due disgiunti e ricoprono
tutto $A$.

\begin{esempio}
\leavevmode
\begin{enumerate}
    \item Se $A \neq \emptyset$, $\mathcal{F} = \{A\}$ è una partizione.
    \item $\mathcal{F} = \{\{x\} \mid x \in A\}$ è una partizione.
    \item $A = \{a, b, c, d, e\}$ e
          $\mathcal{F} = \{\{a, c, e\}, \{b\}, \{d\}\}$ è una partizione
          di $A$.
    \item In $\mathbb{Z}$ siano
          \begin{align*}
          X_0 &= \{3k \mid k \in \mathbb{Z}\} = \{0, \pm 3, \pm 6, \dots\},\\
          X_1 &= \{3k + 1 \mid k \in \mathbb{Z}\}
                = \{1, 4, -2, 7, -5, \dots\},\\
          X_2 &= \{3k + 2 \mid k \in \mathbb{Z}\}
                = \{2, 5, -1, 8, -4, \dots\}.
          \end{align*}
          Allora $\mathcal{F} = \{X_0, X_1, X_2\}$ è una partizione di
          $\mathbb{Z}$.
\end{enumerate}
\end{esempio}

\begin{proposizione}[A]
Sia $\rho$ una relazione di equivalenza su $A$. Allora l'insieme quotiente
$A/_\rho$ è una partizione di $A$.
\end{proposizione}

\begin{proof}
Verifichiamo le proprietà 1, 2 e 3 della definizione di partizione.
\begin{enumerate}
    \item $[a]_\rho \neq \emptyset$: infatti $a \mathrel{\rho} a$, cioè
          $a \in [a]_\rho$.
    \item Se $[a]_\rho \neq [b]_\rho$, vogliamo $[a]_\rho \cap [b]_\rho =
          \emptyset$. Sia, per assurdo, $x \in [a]_\rho \cap [b]_\rho$:
          allora $x \mathrel{\rho} a$ e $x \mathrel{\rho} b$. Per la
          simmetria $a \mathrel{\rho} x$ e $x \mathrel{\rho} b$, e per la
          transitività $a \mathrel{\rho} b$, cioè $a \in [b]_\rho$ e
          $b \in [a]_\rho$. Quindi $[a]_\rho = [b]_\rho$ (hanno gli stessi
          elementi), contro l'ipotesi. Dunque se $[a]_\rho \neq
          [b]_\rho$ non esiste $x \in [a]_\rho \cap [b]_\rho$, cioè
          $[a]_\rho \cap [b]_\rho = \emptyset$.
    \item $\displaystyle \bigcup_{[a]_\rho \in A/_\rho} [a]_\rho = A$:
          se $a \in A$ allora $a \in [a]_\rho$, cioè le classi
          ricoprono tutto $A$. \qedhere
\end{enumerate}
\end{proof}

Viceversa, da una partizione si ricava una relazione di equivalenza.

\begin{definizione}
Sia $\mathcal{F}$ una partizione di $A$. Si definisce la relazione
$R_\mathcal{F}$ su $A$ ponendo
\[
a \mathrel{R_\mathcal{F}} b \iff \exists\, X \in \mathcal{F} \text{ tale
che } a, b \in X.
\]
Cioè $a$ e $b$ sono in relazione se stanno nello stesso pezzo della
partizione.
\end{definizione}

\begin{proposizione}[B]
Se $\mathcal{F}$ è una partizione su $A$, la relazione $R_\mathcal{F}$ è
una relazione di equivalenza su $A$.
\end{proposizione}

\begin{proof}
Verifichiamo che $R_\mathcal{F}$ è riflessiva, simmetrica e transitiva.
\begin{enumerate}
    \item \emph{Riflessiva}: $a \mathrel{R_\mathcal{F}} a$? Poiché
          $\mathcal{F}$ è una partizione, per la proprietà 3
          $\bigcup_{X \in \mathcal{F}} X = A$. Se $a \in A$, allora
          $a \in \bigcup_{X \in \mathcal{F}} X$, cioè esiste
          $X \in \mathcal{F}$ tale che $a \in X$. Quindi
          $a \mathrel{R_\mathcal{F}} a$.
    \item \emph{Simmetrica}: se $a \mathrel{R_\mathcal{F}} b$, cioè esiste
          $X \in \mathcal{F}$ con $a, b \in X$, allora anche
          $b, a \in X$, cioè $b \mathrel{R_\mathcal{F}} a$.
    \item \emph{Transitiva}: siano $a \mathrel{R_\mathcal{F}} b$ e
          $b \mathrel{R_\mathcal{F}} c$. Esistono $X \in \mathcal{F}$ con
          $a, b \in X$ e $Y \in \mathcal{F}$ con $b, c \in Y$. Poiché
          $b \in X \cap Y \neq \emptyset$, ma per la proprietà 2 gli
          elementi della partizione sono disgiunti, deve essere $X = Y$.
          Quindi $a, b, c \in X = Y$, cioè $a, c \in X$, cioè
          $a \mathrel{R_\mathcal{F}} c$. \qedhere
\end{enumerate}
\end{proof}

\begin{teorema}[Fondamentale delle relazioni di equivalenza]
Esiste una corrispondenza 1 a 1 tra l'insieme di tutte le relazioni di
equivalenza su $A$ e l'insieme di tutte le partizioni di $A$.
\end{teorema}

\begin{proof}
Se $\rho$ è una relazione di equivalenza su $A$, allora $A/_\rho$ è una
partizione (Proposizione A). Se $\mathcal{F}$ è una partizione di $A$,
$R_\mathcal{F}$ è una relazione di equivalenza su $A$ (Proposizione B).
\end{proof}

\begin{osservazione}
Le due costruzioni sono l'una l'inversa dell'altra: la classe di $a$
rispetto a $R_\mathcal{F}$ è esattamente l'elemento di $\mathcal{F}$ che
contiene $a$, quindi $A/_{R_\mathcal{F}} = \mathcal{F}$; e due elementi
stanno nella stessa classe di $A/_\rho$ se e solo se sono in relazione
rispetto a $\rho$, quindi $R_{A/_\rho} = \rho$. Questo è ciò che rende la
corrispondenza biunivoca.
\end{osservazione}

\section[Congruenza modulo $n$ (Lezione 3 -- 28 settembre 2026)]{Congruenza modulo $n$}
\noindent\emph{Lezione 3 -- 28 settembre 2026}
\subsection{La congruenza modulo $n$}

\begin{definizione}
Siano $a, b \in \mathbb{Z}$. Si dice che $a$ \emph{divide} $b$, e si scrive
$a \mid b$, se
\[
\exists\, k \in \mathbb{Z} \text{ tale che } b = ka.
\]
Se non esiste un tale $k$ si scrive $a \nmid b$ (``$a$ non divide $b$'').
\end{definizione}

Sia $n \geq 1$ un intero fissato.

\begin{definizione}
La \emph{congruenza modulo $n$} è la relazione $\equiv_n$ su $\mathbb{Z}$
definita da
\[
a \equiv_n b \iff n \mid (b - a) \iff \exists\, k \in \mathbb{Z} \text{
tale che } b - a = kn \quad (b = kn + a).
\]
Si scrive anche $a \equiv b \pmod n$.
\end{definizione}

\begin{proposizione}
$\equiv_n$ è una relazione di equivalenza su $\mathbb{Z}$, e
\[
\mathbb{Z}/_{\equiv_n} = \mathbb{Z}_n = \{[0]_n, [1]_n, \dots, [n-1]_n\}.
\]
\end{proposizione}

\begin{proof}
Verifichiamo che $\equiv_n$ è una relazione di equivalenza.
\begin{enumerate}
    \item \emph{Riflessiva}: $a \equiv_n a$, perché $n \mid (a - a) = 0$,
          cioè $0 = k \cdot n$ con $k = 0$.
    \item \emph{Simmetrica}: se $a \equiv_n b$, allora $n \mid (b - a)$,
          cioè esiste $k \in \mathbb{Z}$ con $b - a = kn$. Posto
          $k' = -k$ si ha $a - b = -kn = k'n$, quindi $n \mid (a - b)$,
          cioè $b \equiv_n a$.
    \item \emph{Transitiva}: siano $a \equiv_n b$ e $b \equiv_n c$.
          Allora $b - a = kn$ e $c - b = hn$ per certi $k, h \in
          \mathbb{Z}$. Sommando,
          \[
          c - a = kn + hn = (k + h)n,
          \]
          quindi $n \mid (c - a)$, cioè $a \equiv_n c$. \qedhere
\end{enumerate}
\end{proof}

La classe di $a$ è
\begin{align*}
[a]_n &= \{b \in \mathbb{Z} \mid a \equiv_n b\}
       = \{b \in \mathbb{Z} \mid n \mid (b - a)\} \\
      &= \{b \in \mathbb{Z} \mid \exists\, k \in \mathbb{Z} \text{ tale
         che } b = kn + a\}
       = \{kn + a \mid k \in \mathbb{Z}\},
\end{align*}
cioè è l'insieme di tutti i numeri che differiscono da $a$ per un multiplo
di $n$. L'insieme quotiente è
$\mathbb{Z}_n = \{[0]_n, [1]_n, \dots, [n-1]_n\}$.

\begin{esempio}
Per $n = 3$ le classi $[0]_3, [1]_3, [2]_3$ sono gli insiemi $X_0, X_1,
X_2$ dell'esempio delle partizioni di $\mathbb{Z}$ (lezione 2): $\mathbb{Z}_3$ è proprio la partizione
di $\mathbb{Z}$ in multipli di $3$, multipli di $3$ più $1$ e multipli di
$3$ più $2$.
\end{esempio}

\section[Relazioni d'ordine (Lezione 3 -- 28 settembre 2026)]{Relazioni d'ordine}
\noindent\emph{Lezione 3 -- 28 settembre 2026}
\subsection{Ordini parziali e totali}

\begin{definizione}
Sia $A$ un insieme e $R$ una relazione su $A$. $R$ è
\emph{antisimmetrica} se
\[
a \mathrel{R} b \ \land\ b \mathrel{R} a \implies a = b.
\]
\end{definizione}

\begin{definizione}
Una relazione $R$ (che si indica con $\preceq$) su $A$ è una
\emph{relazione d'ordine parziale} se è riflessiva, antisimmetrica e
transitiva.
\end{definizione}

\begin{itemize}
    \item $A$ con una relazione d'ordine si dice \emph{parzialmente
          ordinato}.
    \item Si scrive $a \preceq b$; con $a \prec b$ si intende
          $a \preceq b$ e $a \neq b$.
    \item Una relazione d'ordine $\preceq$ è \emph{totale} se per ogni
          $a, b \in A$ vale $a \preceq b$ oppure $b \preceq a$.
\end{itemize}
$A$ con una relazione d'ordine totale si dice \emph{totalmente ordinato}
o \emph{catena}.

\begin{esempio}
\leavevmode
\begin{enumerate}
    \item $(\mathbb{N}, \leq)$ è riflessivo ($n \leq n$), antisimmetrico
          ($n \leq k \land k \leq n \Rightarrow n = k$) e transitivo
          ($n \leq k \land k \leq h \Rightarrow n \leq h$). È una
          relazione d'ordine \textbf{totale}: per ogni $n, k \in
          \mathbb{N}$ vale $n \leq k$ oppure $k \leq n$.

    \item $(\mathbb{Z}, \leq)$ è una relazione d'ordine totale.

    \item Sia $B = \{x, y, z\}$. Allora
          \[
          \mathcal{P}(B) = \big\{\emptyset, \{x\}, \{y\}, \{z\},
          \{x,y\}, \{x,z\}, \{y,z\}, B\big\}.
          \]
          $(\mathcal{P}(B), \subseteq)$ è una relazione d'ordine
          \textbf{parziale}. Non è totale: $\{x, y\}$ e $\{y, z\}$ non
          sono in relazione.
\end{enumerate}
\end{esempio}

\subsection{Il diagramma di Hasse}

\begin{osservazione}
Ad ogni relazione d'ordine si può associare un \emph{diagramma di Hasse}:
\begin{itemize}
    \item i nodi sono esattamente gli elementi dell'insieme;
    \item se $x \preceq y$ (con $x \neq y$) il nodo $x$ è disegnato più in
          basso del nodo $y$;
    \item due nodi $x, y$ sono collegati da un arco se \emph{non} esiste
          $z$ tale che $x \prec z$ e $z \prec y$ (cioè se $y$ ``viene
          subito dopo'' $x$).
\end{itemize}
\end{osservazione}

\begin{center}
\begin{tikzpicture}[scale=1.1, every node/.style={font=\small}]
  % catena N
  \begin{scope}
    \foreach \n/\y in {0/0, 1/1, 2/2, 3/3} {
      \fill (0,\y) circle (1.5pt) node[right=2pt] {$\n$};
    }
    \draw (0,0) -- (0,3.4);
    \node at (0,-0.7) {$(\mathbb{N}, \leq)$};
  \end{scope}

  % P(B)
  \begin{scope}[xshift=6cm]
    \node (e) at (0,0) {$\emptyset$};
    \node (x) at (-2,1.4) {$\{x\}$};
    \node (y) at (0,1.4) {$\{y\}$};
    \node (z) at (2,1.4) {$\{z\}$};
    \node (xy) at (-2,2.9) {$\{x,y\}$};
    \node (xz) at (0,2.9) {$\{x,z\}$};
    \node (yz) at (2,2.9) {$\{y,z\}$};
    \node (b) at (0,4.3) {$B$};
    \foreach \a/\b in {e/x, e/y, e/z, x/xy, x/xz, y/xy, y/yz, z/xz, z/yz,
                       xy/b, xz/b, yz/b}
      \draw (\a) -- (\b);
    \node at (0,-0.9) {$(\mathcal{P}(B), \subseteq)$};
  \end{scope}
\end{tikzpicture}
\end{center}

\begin{esempio}[continua]
\leavevmode
\begin{enumerate}
    \setcounter{enumi}{3}
    \item Sia $X = \{1, 2, 3, 4, 5, 6\}$ con la relazione ``$\mid$''
          (divide). Le coppie in relazione sono
          \begin{center}
          \begin{tabular}{llllll}
          $1 \mid 1$ & $1 \mid 2$ & $1 \mid 3$ & $1 \mid 4$ & $1 \mid 5$ &
          $1 \mid 6$ \\
                     & $2 \mid 2$ & $2 \mid 4$ & $2 \mid 6$ \\
                     & $3 \mid 3$ & $3 \mid 6$ \\
                     & $4 \mid 4$ & $5 \mid 5$ & $6 \mid 6$
          \end{tabular}
          \end{center}
          È una relazione d'ordine parziale, non totale: $4 \nmid 5$ e
          $5 \nmid 4$.

\begin{center}
\begin{tikzpicture}[scale=1.1, every node/.style={font=\small}]
  \node (1) at (1.5,0) {$1$};
  \node (2) at (0,1.4) {$2$};
  \node (3) at (1.5,1.4) {$3$};
  \node (5) at (3,1.4) {$5$};
  \node (4) at (0,2.8) {$4$};
  \node (6) at (1.5,2.8) {$6$};
  \foreach \a/\b in {1/2, 1/3, 1/5, 2/4, 2/6, 3/6}
    \draw (\a) -- (\b);
\end{tikzpicture}
\end{center}

    \item $(\mathbb{N} \setminus \{0\}, \mid)$ è una relazione d'ordine non
          totale.

    \item $\mathbb{Z} \setminus \{0\}$ con ``$\mid$'' (divide) è una
          relazione d'ordine? \textbf{No}: non è antisimmetrica, perché
          $5 \mid -5$ e $-5 \mid 5$ ma $-5 \neq 5$.
\end{enumerate}
\end{esempio}

\begin{osservazione}
La relazione opposta $R^{op}$ scambia le coppie. Se $\preceq$ è una
relazione d'ordine, allora $(\preceq)^{op} = \succeq$ è ancora una relazione
d'ordine: il suo diagramma di Hasse si ottiene capovolgendo quello di
$\preceq$. Ad esempio $(\leq)^{op} = \ \geq$ su $\mathbb{N}$, e
$(\subseteq)^{op} = \ \supseteq$ su $\mathcal{P}(B)$ (con $B$ in basso e
$\emptyset$ in alto).
\end{osservazione}

\subsection{Minimo, massimo, minimale, massimale}

\begin{definizione}
Sia $(A, \preceq)$ un insieme parzialmente ordinato. Un elemento $t \in A$ è:
\begin{itemize}
    \item \emph{minimo} se $\forall\, a \in A \quad t \preceq a$;
    \item \emph{massimo} se $\forall\, a \in A \quad a \preceq t$;
    \item \emph{minimale} se non ci sono elementi minori di $t$, cioè
          $\nexists\, a \in A$ tale che $a \prec t$;
    \item \emph{massimale} se non ci sono elementi maggiori di $t$, cioè
          $\nexists\, a \in A$ tale che $t \prec a$.
\end{itemize}
\end{definizione}

Un minimo è sempre minimale (e un massimo è sempre massimale), ma il
viceversa in generale è falso se l'ordine non è totale.

\begin{esempio}
Negli esempi precedenti:
\begin{itemize}
    \item $(\mathbb{N}, \leq)$: $0$ è minimo e minimale; non esistono
          massimo né massimali.
    \item $(\mathbb{Z}, \leq)$: non esistono minimo, massimo, minimali né
          massimali.
    \item $(\mathcal{P}(B), \subseteq)$: $\emptyset$ è minimo e minimale,
          $B$ è massimo e massimale.
    \item $(\{1, \dots, 6\}, \mid)$: $1$ è minimo e minimale; non esiste
          il massimo; $4$, $5$, $6$ sono massimali.
    \item $(\mathbb{N} \setminus \{0\}, \mid)$: $1$ è minimo e minimale;
          non esistono massimo né massimali.
\end{itemize}
\end{esempio}

\subsection{Maggioranti, minoranti, estremo superiore e inferiore}

\begin{definizione}
Sia $\preceq$ una relazione d'ordine su $A$ e sia $B \subseteq A$ un
sottoinsieme. Un elemento $w \in A$ è:
\begin{itemize}
    \item un \emph{maggiorante} di $B$ se $b \preceq w$ per ogni
          $b \in B$;
    \item un \emph{minorante} di $B$ se $w \preceq b$ per ogni $b \in B$.
\end{itemize}
\end{definizione}

\begin{esempio}
\leavevmode
\begin{enumerate}
    \item In $(\mathbb{N}, \leq)$ sia $B = \{6, 9, 12\}$.
          I minoranti di $B$ sono $\{0, 1, 2, \dots, 6\}$, i maggioranti
          sono $\{12, 13, 14, \dots\}$. Il più grande dei minoranti è
          $6 = \inf(B)$, il più piccolo dei maggioranti è
          $12 = \sup(B)$.

    \item In $(\mathbb{Z}, \leq)$ con $B = \{6, 9, 12\}$: i minoranti sono
          $\{\dots,\allowbreak -1,\allowbreak 0,\allowbreak 1,\allowbreak \dots,\allowbreak 6\}$, i maggioranti
          $\{12, 13, \dots\}$; $6 = \inf(B)$ e $12 = \sup(B)$.

    \item In $(\mathbb{N} \setminus \{0\}, \mid)$ con $B = \{6, 9, 12\}$:
          i minoranti sono i divisori comuni, $\{1, 3\}$; poiché $1 \mid 3$,
          il più grande (rispetto a $\mid$) è $3 = \inf(B)$. I maggioranti
          sono i multipli comuni, $\{36, 72, \dots\}$, cioè i multipli di
          $36$; $36$ divide tutti gli altri, quindi $36 = \sup(B)$.

    \item In $(\mathbb{Q}, \leq)$ sia
          $B = \{x \in \mathbb{Q} \mid \sqrt{2} < x < \sqrt{3}\}$.
          I minoranti sono tutti i razionali $q < \sqrt{2}$, i maggioranti
          tutti i razionali $q > \sqrt{3}$. Tra i minoranti non c'è un
          massimo e tra i maggioranti non c'è un minimo: in $\mathbb{Q}$
          non esistono né $\inf(B)$ né $\sup(B)$.
\end{enumerate}
\end{esempio}

\begin{definizione}
Sia $(A, \preceq)$ parzialmente ordinato e $B \subseteq A$.
\begin{itemize}
    \item Un elemento $t \in A$ è l'\emph{estremo inferiore} di $B$,
          $t = \inf(B)$, se $t$ è il \textbf{massimo} tra i minoranti di
          $B$, cioè
          \[
          t = \inf(B) \iff
          \begin{cases}
            t \preceq b & \forall\, b \in B, \\
            w \preceq b \ \ \forall\, b \in B \implies w \preceq t.
          \end{cases}
          \]
    \item Un elemento $t \in A$ è l'\emph{estremo superiore} di $B$,
          $t = \sup(B)$, se $t$ è il \textbf{minimo} tra i maggioranti di
          $B$, cioè
          \[
          t = \sup(B) \iff
          \begin{cases}
            b \preceq t & \forall\, b \in B, \\
            b \preceq w \ \ \forall\, b \in B \implies t \preceq w.
          \end{cases}
          \]
\end{itemize}
\end{definizione}

\section[Funzioni (Lezione 3 -- 28 settembre 2026)]{Funzioni}
\noindent\emph{Lezione 3 -- 28 settembre 2026}
\subsection{Definizione di funzione}

\begin{definizione}
Siano $A$ e $B$ due insiemi. Una \emph{funzione} $f$ da $A$ a $B$ è un
sottoinsieme del prodotto cartesiano $A \times B$ tale che
\[
\forall\, a \in A \quad \exists!\, b \in B : (a, b) \in f.
\]
\end{definizione}

In altre parole, ogni elemento di $A$ compare come prima componente in
\textbf{una e una sola} coppia di $f$. Si usano le notazioni seguenti:
\begin{itemize}
    \item $f : A \to B$ indica che $f$ è una funzione da $A$ a $B$;
    \item $a \mathrel{f} b$ oppure $f(a) = b$ significa $(a, b) \in f$;
          $b$ si dice \emph{immagine} di $a$;
    \item $A$ è il \emph{dominio} di $f$ e $B$ è il \emph{codominio} di $f$.
\end{itemize}

Il sottoinsieme di $A \times B$
\[
f = \{\, (a, f(a)) \mid a \in A,\ f(a) \in B \,\} \subseteq A \times B
\]
si chiama \emph{grafico} della funzione $f$.

\begin{osservazione}
Secondo la definizione, la funzione \emph{coincide} con il suo grafico.
Le due condizioni che la definiscono sono:
\begin{enumerate}
    \item \textbf{esistenza}: ogni $a \in A$ ha almeno un'immagine in $B$
          (nessun elemento del dominio resta ``senza freccia'');
    \item \textbf{unicità}: ogni $a \in A$ ha \emph{una sola} immagine
          (da ogni elemento del dominio parte una sola freccia).
\end{enumerate}
Invece \emph{non} si chiede che ogni elemento di $B$ sia immagine di
qualcuno, né che sia immagine di uno solo: nell'esempio~2 qui sotto $a$ e
$b$ hanno ciascuno due elementi che arrivano su di loro.
\end{osservazione}

\subsection{Esempi}

\begin{esempio}
\leavevmode
\begin{enumerate}
    \item Siano $A = \{x, y, z\}$ e $B = \{1, 2, 3, 4\}$ e sia
          $f : A \to B$ data da
          \[
          f(x) = 1, \qquad f(y) = 3, \qquad f(z) = 4.
          \]
          Quindi $x \mathrel{f} 1$, cioè $(x, 1) \in f$, e il grafico è
          $f = \{(x,1), (y,3), (z,4)\}$. L'elemento $2 \in B$ non è
          immagine di nessun elemento di $A$.

    \item Siano $B = \{1, 2, 3, 4\}$ e $C = \{a, b\}$ e sia
          $g : B \to C$ data da
          \[
          g(1) = b, \quad g(2) = b, \quad g(3) = a, \quad g(4) = a.
          \]
\end{enumerate}

\begin{center}
\begin{tikzpicture}[scale=0.9, every node/.style={font=\small}]
  % funzione f : A -> B
  \draw (0,0) ellipse (0.9 and 1.4);
  \draw (3,0) ellipse (0.9 and 1.7);
  \node at (0,1.75) {$A$};
  \node at (3,2.05) {$B$};
  \node at (1.5,2.0) {$f$};
  \foreach \n/\y in {x/0.8, y/0, z/-0.8} {
    \fill (-0.3,\y) circle (1.5pt) node[left] {$\n$};
  }
  \foreach \n/\y in {1/1.1, 2/0.35, 3/-0.4, 4/-1.15} {
    \fill (2.7,\y) circle (1.5pt) node[right] {$\n$};
  }
  \draw[->] (-0.3,0.8) -- (2.7,1.1);
  \draw[->] (-0.3,0) -- (2.7,-0.4);
  \draw[->] (-0.3,-0.8) -- (2.7,-1.15);

  % funzione g : B -> C
  \begin{scope}[xshift=8cm]
    \draw (0,0) ellipse (0.9 and 1.7);
    \draw (3,0) ellipse (0.9 and 1.0);
    \node at (0,2.05) {$B$};
    \node at (3,1.35) {$C$};
    \node at (1.5,2.0) {$g$};
    \foreach \n/\y in {1/1.1, 2/0.35, 3/-0.4, 4/-1.15} {
      \fill (-0.3,\y) circle (1.5pt) node[left] {$\n$};
    }
    \fill (2.7,0.5) circle (1.5pt) node[right] {$b$};
    \fill (2.7,-0.5) circle (1.5pt) node[right] {$a$};
    \draw[->] (-0.3,1.1) -- (2.7,0.5);
    \draw[->] (-0.3,0.35) -- (2.7,0.5);
    \draw[->] (-0.3,-0.4) -- (2.7,-0.5);
    \draw[->] (-0.3,-1.15) -- (2.7,-0.5);
  \end{scope}
\end{tikzpicture}
\end{center}

\begin{enumerate}
    \setcounter{enumi}{2}
    \item $f : \mathbb{N} \to \mathbb{N}$, \quad $f(n) = 2n$.
          Si scrive anche $n \mapsto f(n) = 2n$.

    \item $f : \mathbb{Z} \to \mathbb{N}$, \quad
          $k \mapsto |k| =
          \begin{cases}
            k & \text{se } k \geq 0, \\
            -k & \text{se } k < 0.
          \end{cases}$

    \item $h : \mathbb{R} \to \mathbb{R}$, \quad $h(x) = x^2 - 2$.

    \item $L : \mathbb{R} \to \mathbb{R}$, \quad $L(x) = 5x + 3$.

    \item $F : \mathbb{Z} \to \mathbb{Z}$, \quad $F(k) = 3$.

    \item $F : \mathbb{Z} \to \{3\}$, \quad $F(k) = 3$.
\end{enumerate}
\end{esempio}

Gli ultimi due esempi sono funzioni \emph{costanti}: tutti gli elementi del
dominio hanno la stessa immagine. Notiamo che la legge $F(k) = 3$ è la
stessa, ma le due funzioni hanno codomini diversi ($\mathbb{Z}$ nel primo
caso, $\{3\}$ nel secondo).

\begin{definizione}
Una funzione $f : A \to B$ tale che
\[
f(x_1) = f(x_2) \qquad \forall\, x_1, x_2 \in A
\]
si dice \emph{funzione costante}.
\end{definizione}

\section[Funzioni: la relazione associata (Lezione 4 -- 29 settembre 2026)]{Funzioni: la relazione associata}
\noindent\emph{Lezione 4 -- 29 settembre 2026}
Si riprendono le funzioni introdotte nella lezione 3 (definizione, dominio, codominio, grafico, funzioni costanti).

\subsection{La relazione associata a una funzione}

\begin{definizione}
Sia $f : A \to B$ una funzione. La \emph{relazione associata a $f$}, indicata
con $R_f$, è la relazione su $A$ definita da
\[
a_1 \mathrel{R_f} a_2 \iff f(a_1) = f(a_2).
\]
\end{definizione}

Due elementi di $A$ sono quindi in relazione se $f$ associa loro la stessa
immagine.

\begin{esempio}
Siano $A = \{1, 2, 3, 4\}$ e $B = \{x, y, z\}$ e sia $f : A \to B$ data da
\[
f(1) = f(2) = f(3) = x, \qquad f(4) = y.
\]
Allora $1 \mathrel{R_f} 2$, $3 \mathrel{R_f} 3$, $4 \mathrel{R_f} 4$, \dots{}
Le classi di equivalenza sono
\[
[1]_{R_f} = \{1, 2, 3\} = [2]_{R_f} = [3]_{R_f}, \qquad
[4]_{R_f} = \{4\},
\]
e l'insieme quotiente è $A/_{R_f} = \{[1]_{R_f}, [4]_{R_f}\}$.
\end{esempio}

\begin{teorema}
Sia $f : A \to B$ una funzione. La relazione $R_f$ è una relazione di
equivalenza su $A$.
\end{teorema}

\begin{proof}
Lasciata per esercizio a lezione; la verifica è
immediata, perché ``avere la stessa immagine'' eredita le proprietà
dell'uguaglianza in $B$.
\begin{enumerate}
    \item \emph{Riflessiva}: per ogni $a \in A$ vale $f(a) = f(a)$, quindi
          $a \mathrel{R_f} a$.
    \item \emph{Simmetrica}: se $a_1 \mathrel{R_f} a_2$, cioè
          $f(a_1) = f(a_2)$, allora $f(a_2) = f(a_1)$, cioè
          $a_2 \mathrel{R_f} a_1$.
    \item \emph{Transitiva}: se $a_1 \mathrel{R_f} a_2$ e
          $a_2 \mathrel{R_f} a_3$, allora $f(a_1) = f(a_2)$ e
          $f(a_2) = f(a_3)$, quindi $f(a_1) = f(a_3)$, cioè
          $a_1 \mathrel{R_f} a_3$. \qedhere
\end{enumerate}
\end{proof}

\section[Principio di induzione (Lezione 4 -- 29 settembre 2026)]{Principio di induzione}
\noindent\emph{Lezione 4 -- 29 settembre 2026}
\subsection{Il principio di induzione}

\begin{osservazione}
Sia $U \subseteq \mathbb{N}$ tale che
\begin{enumerate}
    \item $0 \in U$;
    \item $k \in U \Rightarrow k + 1 \in U$.
\end{enumerate}
Allora $U = \mathbb{N}$.
\end{osservazione}

Infatti $0 \in U$, quindi $1 \in U$, quindi $2 \in U$, e così via: ``si sale
a gradini'' e si raggiungono tutti i naturali. Se $p(n)$ è un predicato
sui naturali, si applica questa idea all'insieme $U$ dei $n$ per cui
$p(n)$ è vero.

\begin{teorema}[Principio di induzione]
Sia $p(n)$ un predicato. Supponiamo che:
\begin{enumerate}
    \item $p(0)$ sia vero \hfill [\emph{base dell'induzione}];
    \item $\forall\, k \ \big(p(k) \text{ vero} \Rightarrow p(k+1)
          \text{ vero}\big)$ \hfill [\emph{passo induttivo}].
\end{enumerate}
Allora $\forall\, n \in \mathbb{N}$ $p(n)$ è vero.
\end{teorema}

Nel passo induttivo l'ipotesi $p(k)$ si chiama \emph{ipotesi induttiva} e
$p(k+1)$ è la \emph{tesi da dimostrare}.

\begin{esempio}
Dimostrare che per ogni $n$ naturale $2 \mid n(n+1)$.

Sia $p(n)$ il predicato ``$2 \mid n(n+1)$''.

\textbf{1) Caso base.} $p(0)$: $2 \mid 0 \cdot (0 + 1) = 0$. Infatti
$2 \mid 0$ significa $0 = k \cdot 2$ con $k \in \mathbb{Z}$, ed è vero per
$k = 0$: $0 = 0 \cdot 2$.

\textbf{2) Passo induttivo.} Dobbiamo dimostrare $p(k) \Rightarrow
p(k+1)$.
\begin{itemize}
    \item \emph{Ipotesi induttiva} $p(k)$: $2 \mid k(k+1)$.
    \item \emph{Tesi da dimostrare} $p(k+1)$: $2 \mid (k+1)\big((k+1)+1\big)$.
\end{itemize}
Calcoliamo
\[
(k+1)(k+2) = (k+1)k + 2(k+1).
\]
Per l'ipotesi induttiva $2 \mid k(k+1)$, cioè $k(k+1) = 2m$ con
$m \in \mathbb{Z}$. Quindi
\[
(k+1)(k+2) = 2m + 2(k+1) = 2\,(m + k + 1),
\]
cioè $(k+1)(k+2)$ è multiplo di $2$, ovvero $2 \mid (k+1)(k+2)$, che è la
tesi da dimostrare.
\end{esempio}

\subsection{Induzione a partire da $\bar n$}

Se il predicato è vero solo da un certo naturale in poi, si parte da
quel punto.

\begin{teorema}
Sia $p(n)$ una proposizione e sia $\bar n \in \mathbb{N}$. Supponiamo che:
\begin{enumerate}
    \item $p(\bar n)$ sia vero;
    \item $\forall\, k \geq \bar n \ \big(p(k) \text{ vero} \Rightarrow
          p(k+1) \text{ vero}\big)$.
\end{enumerate}
Allora $\forall\, n \geq \bar n$ $p(n)$ è vera.
\end{teorema}

\begin{esempio}
Il predicato $p(n)$: ``$n - 5 > 0$'' è vero da $6$ in poi: si applica il
teorema con $\bar n = 6$.
\end{esempio}

\begin{esempio}
Dimostrare che per ogni $n \geq 1$
\[
1 + 2 + \dots + n = \frac{n(n+1)}{2}. \tag*{$p(n)$}
\]

\textbf{1) Caso base} $\bar n = 1$. $p(1)$: $1 = \dfrac{1 \cdot (1+1)}{2}$,
vera.

\textbf{2) Passo induttivo.}
\begin{itemize}
    \item \emph{Ipotesi induttiva} $p(k)$:
          $1 + \dots + k = \dfrac{k(k+1)}{2}$.
    \item \emph{Tesi da dimostrare} $p(k+1)$:
          $1 + \dots + k + (k+1) = \dfrac{(k+1)(k+2)}{2}$.
\end{itemize}
Si ha
\begin{align*}
1 + 2 + \dots + k + (k+1) &= (1 + \dots + k) + (k+1)\\
&= \frac{k(k+1)}{2} + (k+1) \qquad \text{(per ipotesi induttiva)}\\
&= \frac{k(k+1) + 2(k+1)}{2} = \frac{(k+1)(k+2)}{2},
\end{align*}
che è la tesi. Per induzione $p(n)$ vale per ogni $n \geq 1$.
\end{esempio}

\section[Successioni (Lezione 4 -- 29 settembre 2026)]{Successioni}
\noindent\emph{Lezione 4 -- 29 settembre 2026}
\subsection{Successioni}

\begin{definizione}
Una \emph{successione} è una funzione $f : \mathbb{N} \to X$ (oppure con
dominio $\mathbb{N} \setminus \{0\}$), dove $X$ è un insieme qualsiasi.
\end{definizione}

L'immagine $f(n)$ di $n$ si indica con $x_n \in X$ e si chiama anche
l'\emph{$n$-esimo termine} della successione. L'insieme immagine di una
successione è
\[
\operatorname{Im} f = f(\mathbb{N}) = \{x_n\}_{n \in \mathbb{N}}
\subseteq X.
\]

\begin{esempio}
\leavevmode
\begin{enumerate}
    \item $a : \mathbb{N} \setminus \{0\} \to \mathbb{Q}$,
          $n \mapsto \dfrac{1}{n}$. Si scrive $a_n = \dfrac{1}{n} = a(n)$
          e l'insieme dei termini è
          $\{a_n\} = \left\{\dfrac{1}{n}\right\}_{n \in \mathbb{N}
          \setminus \{0\}}$.

    \item $b : \mathbb{N} \to \mathbb{N}$, $n \mapsto 2n + 3$. I primi
          termini sono $b_0 = 2 \cdot 0 + 3 = 3$,
          $b_1 = 2 \cdot 1 + 3 = 5$, \dots

    \item $c : \mathbb{N} \to \mathbb{N}$, $n \mapsto 2^n$. Ad esempio
          $c_3 = 2^3 = 8$, \dots
\end{enumerate}
\end{esempio}

\section[Sommatorie, ricorsione e polinomi (Lezione 5 -- 5 ottobre 2026)]{Sommatorie, ricorsione e polinomi}
\noindent\emph{Lezione 5 -- 5 ottobre 2026}
\subsection{Sommatorie}

Sia $f \colon \mathbb{N} \to \mathbb{R}$ una successione, $f(n) = a_n$. Ad
esempio, per $a_n = n$ si ha $a_1 = 1$, $a_2 = 2$, \dots, $a_{100} = 100$.

Sommare i primi due termini significa calcolare $a_1 + a_2 = 1 + 2$; i primi
cinque, $a_1 + a_2 + a_3 + a_4 + a_5 = 1 + 2 + 3 + 4 + 5$. Per scrivere somme
con molti termini si usa il simbolo di \emph{sommatoria}:
\[
\sum_{i=1}^{n} f(i) = f(1) + f(2) + f(3) + \dots + f(n).
\]

\begin{esempio}
La somma dei primi $100$ numeri è
\[
1 + 2 + \dots + 100 = \sum_{i=1}^{100} a_i = \sum_{i=1}^{100} i .
\]
\end{esempio}

\begin{definizione}
Sia $a_n \colon \mathbb{N} \to \mathbb{R}$ una successione. Le si associa
un'altra successione
\[
s_k = a_1 + a_2 + \dots + a_k = \sum_{i=1}^{k} a_i .
\]
La successione $\{s_k\}$ si dice \emph{serie associata} alla successione
$\{a_n\}$.
\end{definizione}

\subsubsection{Proprietà della sommatoria}

\begin{proposizione}
Siano $\{a_n\}$, $\{b_n\}$ due successioni in $\mathbb{R}$ e $\bar n$ l'indice
iniziale della somma. Si hanno le seguenti proprietà:
\begin{enumerate}
    \item $\displaystyle \sum_{i=\bar n}^{n} (\lambda a_i)
          = \lambda \left( \sum_{i=\bar n}^{n} a_i \right)$, \quad
          $\lambda \in \mathbb{R}$;
    \item $\displaystyle \sum_{i=\bar n}^{n} (a_i + b_i)
          = \left( \sum_{i=\bar n}^{n} a_i \right)
          + \left( \sum_{i=\bar n}^{n} b_i \right)$;
    \item $\displaystyle \sum_{i=\bar n}^{n} a_i
          = \sum_{i=\bar n}^{k} a_i + \sum_{i=k+1}^{n} a_i$, \quad
          con $\bar n \leq k < n$;
    \item (cambio d'indice)
          $\displaystyle \sum_{i=1}^{n} a_i = \sum_{j=1}^{n} a_j
          = \sum_{k=1}^{n} a_k = \sum_{k=0}^{n-1} a_{k+1}$.
\end{enumerate}
\end{proposizione}

\subsection{Ricorsione}

\begin{definizione}
Una successione $\{a_n\}_{n \in \mathbb{N}}$ di numeri reali è
\emph{ricorsiva} se è definita specificando:
\begin{enumerate}
    \item il valore dei primi $m$ termini $a_1, a_2, \dots, a_m$;
    \item una funzione $f \colon \mathbb{R}^m \to \mathbb{R}$ tale che
          \[
            a_i = f(a_{i-1}, a_{i-2}, \dots, a_{i-m}) \qquad \forall\, i > m .
          \]
\end{enumerate}
\end{definizione}

\begin{esempio}
\begin{enumerate}
    \item Sia $a_0 = 0$ e $a_n = 1 + a_{n-1}$ per $n \geq 1$. Allora
          $a_1 = 1 + a_0 = 1$, $a_2 = 1 + a_1 = 2$, $a_3 = 1 + a_2 = 3$, \dots,
          $a_n = n$ (formula chiusa).
    \item \emph{Successione di Fibonacci}: $a_1 = 1$, $a_2 = 1$,
          $a_n = a_{n-1} + a_{n-2}$. Si ottiene $1, 1, 2, 3, 5, 8, \dots$
          (ad esempio $a_3 = a_2 + a_1 = 1 + 1 = 2$, $a_4 = a_3 + a_2 = 3$,
          $a_5 = a_4 + a_3 = 5$, $a_6 = a_5 + a_4 = 8$).
    \item Data $\{a_n\}$, la sua serie associata $\{s_n\}$ si definisce per
          ricorsione:
          \[
            \begin{cases}
              s_0 = a_0\\
              s_n = s_{n-1} + a_n
            \end{cases}
          \]
    \item $a_1 = 5$ e $a_n = 8$ per ogni $n > 1$: tutti i termini dopo il primo
          valgono $8$.
    \item \emph{Fattoriale}: $a_1 = 1$, $a_n = n \cdot a_{n-1}$. Allora
          $a_2 = 2 \cdot a_1 = 2$, $a_3 = 3 \cdot a_2 = 3 \cdot 2 \cdot 1$,
          $a_4 = 4 \cdot 3 \cdot 2 \cdot 1$, \dots, e in generale
          $a_n = n \cdot (n-1) \cdot (n-2) \cdots 3 \cdot 2 \cdot 1$.
          Si scrive $a_n = n!$ (formula esplicita), con $1! = 1$ e per
          convenzione $0! = 1$.
\end{enumerate}
\end{esempio}

\begin{definizione}
Data una successione $\{a_n\}$, una \emph{formula esplicita} o
\emph{formula chiusa} per $\{a_n\}$ è una formula $p(n)$ che esprime
direttamente il valore della successione $a_n$ a partire da $n$.

Se $\{a_n\}$ è definita per ricorsione, una formula chiusa per $\{a_n\}$ è
un'espressione $p(n)$ che verifica le condizioni iniziali e verifica la
relazione ricorsiva.
\end{definizione}

\begin{esempio}
Sia
\[
  \begin{cases}
    a_0 = 2\\
    a_n = a_{n-1} + 5
  \end{cases}
\]
Si ottiene $a_1 = 2 + 5 = 7$, $a_2 = 12$, $a_3 = 17$, $a_4 = 22$. La formula
chiusa $a_n = 5n + 2$ è quella giusta? Verifichiamo.
\begin{enumerate}
    \item \emph{Caso base} ($n = 0$): la formula dà $0 \cdot 5 + 2 = 2$, che
          coincide con il dato iniziale $a_0 = 2$.
    \item \emph{Relazione ricorsiva} $a_n = a_{n-1} + 5$: sostituendo la
          formula con $n-1$ al posto di $n$,
          \[
            a_{n-1} + 5 = \bigl( (n-1) \cdot 5 + 2 \bigr) + 5
            = 5n - 5 + 2 + 5 = 5n + 2 = a_n .
          \]
\end{enumerate}
Sono verificate entrambe: la formula chiusa è verificata.
\end{esempio}

\begin{esercizio}
Verificare che per
$\begin{cases} a_1 = 1\\ a_n = n + a_{n-1} \end{cases}$
la formula chiusa è $a_n = \dfrac{n(n+1)}{2}$.
\end{esercizio}

\emph{Svolgimento (aggiunto).} Per $n = 1$ la formula dà $\frac{1 \cdot 2}{2} = 1 = a_1$.
Inoltre
$a_{n-1} + n = \frac{(n-1)n}{2} + n = \frac{n^2 - n + 2n}{2} = \frac{n(n+1)}{2}$.

\subsubsection{Progressioni aritmetiche}

\begin{definizione}
Una successione $a_n \colon \mathbb{N} \to \mathbb{R}$ è una
\emph{progressione aritmetica} se
\[
  a_n - a_{n-1} = d, \qquad d \in \mathbb{R}.
\]
Definita in maniera ricorsiva, $a_0 = a$ e $a_n = d + a_{n-1}$, ha formula
chiusa
\[
  a_n = a + n \cdot d .
\]
\end{definizione}

\begin{esercizio}
Data $\begin{cases} a_0 = 8\\ a_n = a_{n-1} - 3 \end{cases}$, trovare la formula
chiusa e verificarla.
\end{esercizio}

\emph{Svolgimento (aggiunto).} È una progressione aritmetica con $a = 8$ e
$d = -3$, quindi $a_n = 8 - 3n$: per $n = 0$ vale $8$ e
$a_{n-1} - 3 = 8 - 3(n-1) - 3 = 8 - 3n = a_n$.

\subsection{Polinomi}

\begin{definizione}
Un \emph{polinomio nell'incognita $t$ a coefficienti in $\mathbb{R}$} è una
funzione $p \colon \mathbb{R} \to \mathbb{R}$ del tipo
\[
  p(t) = a_0 + a_1 t + a_2 t^2 + a_3 t^3 + \dots + a_n t^n
       = \sum_{i=0}^{n} a_i t^i ,
\]
dove $a_0, a_1, \dots, a_n \in \mathbb{R}$ sono i \emph{coefficienti} e $a_0$ è
il \emph{termine noto}.

Il più grande indice $k \in \mathbb{N}$ tale che $a_i = 0$ per $i > k$ e
$a_k \neq 0$ è il \emph{grado} del polinomio, $\operatorname{gr}(p) = k$. Se
$k$ non esiste, $\operatorname{gr}(p) = -\infty$.
\end{definizione}

\begin{esempio}
\begin{itemize}
    \item $o(t) = 0$ è il \emph{polinomio nullo}, con
          $\operatorname{gr}(o) = -\infty$.
    \item $q(t) = 3$ è un \emph{polinomio costante}
          ($a_0 = 3$, $a_1 = a_2 = \dots = 0$) e $\operatorname{gr}(q) = 0$.
    \item $p(t) = 5 - t + t^2$ ha $\operatorname{gr}(p) = 2$.
\end{itemize}
\end{esempio}

Si indica con $\mathbb{R}[t]$ l'insieme dei polinomi nell'incognita $t$ a
coefficienti reali, e con
\[
  \mathbb{R}_d[t] = \{ p \in \mathbb{R}[t] \mid \operatorname{gr}(p) \leq d \}
\]
quelli di grado al più $d$. Ad esempio $5 + t^2$ e $1 - 3t - t^2$ stanno in
$\mathbb{R}_2[t]$, il polinomio nullo $o$ sta in $\mathbb{R}_0[t]$, mentre
$5 + t^2 \notin \mathbb{R}_1[t]$.

\begin{osservazione}
Due polinomi $p = \sum_{i=0}^{n} a_i t^i$ e $q = \sum_{i=0}^{n} b_i t^i$ sono
uguali, $p = q$, \textbf{se e solo se} $a_i = b_i$ per ogni $i$. Ad esempio
$(x+2)^2 = x^2 + 4x + 4$.
\end{osservazione}

\subsubsection{Operazioni}

Dati $p = \sum_{i=0}^{n} a_i t^i$, $q = \sum_{i=0}^{n} b_i t^i$ e
$\lambda \in \mathbb{R}$, si definiscono:
\begin{itemize}
    \item il \emph{prodotto per uno scalare}:
          $(\lambda p) = \sum_{i=0}^{n} (\lambda a_i)\, t^i$;
    \item la \emph{somma}: $(p + q) = \sum_{i=0}^{n} (a_i + b_i)\, t^i$;
    \item il polinomio $t p = \sum_{i=0}^{n} a_i t^{i+1}
          = t a_0 + a_1 t^2 + \dots + a_n t^{n+1}
          = \sum_{i=1}^{n+1} a_{i-1} t^i$;
    \item il \emph{prodotto}:
          $p \cdot q = (a_0 + \dots + a_n t^n) \cdot q
          = a_0 q + a_1 t q + \dots + a_n t^n q$.
\end{itemize}

\begin{esempio}
Siano $p = 6 + 3t - t^2$ e $q = -6 + 5t$. Allora
\begin{align*}
p + q &= (6 - 6) + (3 + 5)\,t + (-1 + 0)\,t^2 = 8t - t^2,\\
p \cdot q &= (6 + 3t - t^2)(-6 + 5t)
  = 6(-6 + 5t) + 3t(-6 + 5t) - t^2(-6 + 5t)\\
  &= -36 + 30t - 18t + 15t^2 + 6t^2 - 5t^3
  = -36 + 12t + 21t^2 - 5t^3 .
\end{align*}
\end{esempio}

\begin{osservazione}
Per i gradi valgono
\[
  \operatorname{gr}(p + q) \leq \max\{\operatorname{gr}(p), \operatorname{gr}(q)\},
  \qquad
  \operatorname{gr}(p \cdot q) = \operatorname{gr}(p) + \operatorname{gr}(q).
\]
\end{osservazione}

\subsection{Divisione tra polinomi}

\begin{teorema}[Divisione con resto]
Siano $f, g \in \mathbb{R}[t]$ con $g \neq 0$. Allora esistono e sono unici due
polinomi $q, r \in \mathbb{R}[t]$ tali che
\[
  f = q \cdot g + r , \qquad \operatorname{gr}(r) < \operatorname{gr}(g).
\]
$q$ è il \emph{quoziente} e $r$ il \emph{resto}.
\end{teorema}

\begin{esempio}
Siano $f(t) = t^4 - t^3 + t - 1$ e $g(t) = t^2 - t$. Si ha
$t^2 \cdot g = t^4 - t^3$, quindi
\[
  f - t^2 g = t - 1 ,
\]
che ha grado $1 < \operatorname{gr}(g) = 2$: la divisione si ferma. Dunque
$q = t^2$, $r = t - 1$ e
\[
  f(t) = t^2 \cdot g + (t - 1).
\]
\end{esempio}

\begin{esempio}
Siano $f(t) = t^4 - t^3 + t - 1$ e $g(t) = t + 1$, con
$\operatorname{gr}(g) = 1$. Ad ogni passo si sottrae al resto parziale il
multiplo di $g$ che ne elimina il termine di grado massimo:
\[
\begin{array}{rcl}
  t^4 - t^3 + 0t^2 + t - 1 \;-\; t^3 g &=& -2t^3 + 0t^2 + t - 1 \\
  -2t^3 + 0t^2 + t - 1 \;-\; (-2t^2)\, g &=& 2t^2 + t - 1 \\
  2t^2 + t - 1 \;-\; (2t)\, g &=& -t - 1 \\
  -t - 1 \;-\; (-1)\, g &=& 0
\end{array}
\]
Il quoziente è $q = t^3 - 2t^2 + 2t - 1$ e il resto è $r = 0$
($\operatorname{gr}(0) = -\infty < \operatorname{gr}(g)$). Quindi
\[
  f(t) = (t^3 - 2t^2 + 2t - 1)\, g + 0 .
\]
\end{esempio}

\begin{definizione}
Siano $f, g \in \mathbb{R}[t]$ con $g \neq 0$. Si dice che $g$ \emph{divide}
$f$, e si scrive $g \mid f$, se esiste $q \in \mathbb{R}[t]$ tale che
$f = q \cdot g$ (cioè se il resto della divisione è nullo).
\end{definizione}

\section[Radici di polinomi e ricorsione lineare (Lezione 6 -- 6 ottobre 2026)]{Radici di polinomi e ricorsione lineare}
\noindent\emph{Lezione 6 -- 6 ottobre 2026}
\subsection{Radici di polinomi}

\begin{definizione}
Sia $f \in \mathbb{R}[t]$. Un numero $\alpha \in \mathbb{R}$ si dice \emph{radice}
di $f$ se $f(\alpha) = 0$.
\end{definizione}

Il legame con la divisibilità, vista nella lezione precedente, è il seguente.

\begin{teorema}[del fattore]
Sia $f \in \mathbb{R}[t]$ e sia $\alpha \in \mathbb{R}$. Allora
\[
  \alpha \text{ è una radice di } f
  \iff
  (t - \alpha) \mid f .
\]
\end{teorema}

\begin{proof}
Dividiamo $f$ per $t - \alpha$ (che ha grado $1$): per il teorema di divisione
con resto esistono $q, r \in \mathbb{R}[t]$ tali che
\[
  f(t) = q(t)\,(t - \alpha) + r(t), \qquad \operatorname{gr}(r) < 1 ,
\]
quindi $r$ è un polinomio costante. Valutando in $t = \alpha$ si ottiene
\[
  f(\alpha) = q(\alpha)\cdot 0 + r = r .
\]
Dunque $f(\alpha) = 0$ se e solo se $r = 0$, cioè se e solo se $t - \alpha$
divide $f$.
\end{proof}

\begin{osservazione}
Il resto della divisione di $f$ per $t - \alpha$ è quindi esattamente il valore
$f(\alpha)$. Se $\alpha$ è radice di $f$ si può scrivere $f = (t - \alpha)\, q$,
con $\operatorname{gr}(q) = \operatorname{gr}(f) - 1$: trovare una radice
permette di abbassare il grado del problema.
\end{osservazione}

\begin{definizione}
Sia $f \neq 0$. Una radice $\alpha$ ha \emph{molteplicità} (algebrica)
$m \geq 1$ se $(t - \alpha)^m \mid f$ ma $(t - \alpha)^{m+1} \nmid f$, cioè se
$f = (t - \alpha)^m\, g$ con $g(\alpha) \neq 0$. Le radici di molteplicità $1$
si dicono \emph{semplici}.
\end{definizione}

\subsubsection{Ricerca delle radici e scomposizione}

\begin{osservazione}
Se $f = a_n t^n + \dots + a_1 t + a_0$ ha coefficienti \emph{interi} e
$\frac{p}{q}$ (frazione ridotta) è una radice razionale di $f$, allora
$p \mid a_0$ e $q \mid a_n$. In particolare, se $f$ è monico, le radici
razionali sono intere e dividono $a_0$: basta provare i divisori di $a_0$
(con entrambi i segni).
\end{osservazione}

\begin{esempio}[radici semplici]
Scomponiamo $f(t) = t^3 - 2t^2 - 5t + 6$. Le possibili radici intere sono i
divisori di $6$: $\pm 1, \pm 2, \pm 3, \pm 6$. Si trova
\[
  f(1) = 1 - 2 - 5 + 6 = 0 ,
\]
quindi $1$ è radice e $(t - 1) \mid f$. Per trovare il quoziente si può
usare la regola di Ruffini: si riportano i coefficienti di $f$ e si procede
da sinistra a destra, moltiplicando per $\alpha = 1$ e sommando.
\[
\begin{array}{r|rrrr}
   & 1 & -2 & -5 & 6 \\
 1 &   &  1 & -1 & -6 \\ \hline
   & 1 & -1 & -6 & \;0
\end{array}
\]
L'ultimo numero è il resto ($=f(1)=0$), gli altri sono i coefficienti del
quoziente: $f = (t - 1)(t^2 - t - 6)$. Il fattore di secondo grado si
scompone cercando due numeri con somma $-1$ e prodotto $-6$:
$t^2 - t - 6 = (t - 3)(t + 2)$. Quindi
\[
  f(t) = (t - 1)(t - 3)(t + 2),
\]
e le radici sono $1$, $3$, $-2$, tutte semplici.
\end{esempio}

\begin{esempio}[radice multipla]
Sia $f(t) = t^3 - 3t^2 + 4$. I divisori di $4$ sono $\pm 1, \pm 2, \pm 4$ e
\[
  f(2) = 8 - 12 + 4 = 0 , \qquad f(-1) = -1 - 3 + 4 = 0 .
\]
Dividendo per $t - 2$ (Ruffini con $\alpha = 2$ sui coefficienti
$1, -3, 0, 4$) si ottiene $f = (t - 2)(t^2 - t - 2)$ e poi
$t^2 - t - 2 = (t - 2)(t + 1)$. Dunque
\[
  f(t) = (t - 2)^2 (t + 1) .
\]
La radice $2$ ha molteplicità $2$, la radice $-1$ ha molteplicità $1$.
\end{esempio}

\subsection{Ricorsione lineare}

\begin{definizione}
Una successione $(a_n)_{n \geq 0}$ è \emph{ricorsiva lineare omogenea di grado
$k$} (a coefficienti costanti) se esistono $c_1, \dots, c_k \in \mathbb{R}$,
con $c_k \neq 0$, tali che per ogni $n \geq k$
\[
  a_n = c_1\, a_{n-1} + c_2\, a_{n-2} + \dots + c_k\, a_{n-k} ,
\]
e se sono assegnati i primi $k$ valori $a_0, \dots, a_{k-1}$ (\emph{condizioni
iniziali}). Il polinomio
\[
  p(t) = t^k - c_1\, t^{k-1} - c_2\, t^{k-2} - \dots - c_{k-1}\, t - c_k
\]
si dice \emph{polinomio caratteristico} della ricorrenza.
\end{definizione}

``Lineare'' significa che $a_n$ è combinazione lineare dei termini precedenti,
``omogenea'' che non compaiono termini aggiuntivi che non dipendono dalla
successione.

\begin{osservazione}[perché il polinomio caratteristico]
Cerchiamo soluzioni della forma $a_n = \alpha^n$ con $\alpha \neq 0$.
Sostituendo nella ricorrenza e dividendo per $\alpha^{n-k}$ si ottiene
\[
  \alpha^k = c_1\, \alpha^{k-1} + \dots + c_k
  \iff p(\alpha) = 0 .
\]
Quindi $\alpha^n$ risolve la ricorrenza se e solo se $\alpha$ è una radice del
polinomio caratteristico. I teoremi seguenti dicono come combinare queste
soluzioni per soddisfare anche le condizioni iniziali.
\end{osservazione}

\begin{esempio}
\begin{itemize}
  \item $a_n = 3\, a_{n-1}$ ha grado $1$ e $p(t) = t - 3$; con $a_0$ dato si ha
        $a_n = a_0\, 3^n$.
  \item Fibonacci: $F_n = F_{n-1} + F_{n-2}$, con $F_0 = 0$, $F_1 = 1$. Grado
        $2$, $p(t) = t^2 - t - 1$.
  \item $a_n = 5\, a_{n-1} - 6\, a_{n-2}$ ha grado $2$ e
        $p(t) = t^2 - 5t + 6$.
  \item $a_n = 6 a_{n-1} - 11 a_{n-2} + 6 a_{n-3}$ ha grado $3$ e
        $p(t) = t^3 - 6t^2 + 11t - 6$.
\end{itemize}
\end{esempio}

\subsubsection{Teorema A: radici distinte}

\begin{teorema}[A]
Sia $(a_n)$ una successione ricorsiva lineare omogenea di grado $k$ il cui
polinomio caratteristico $p$ ha $k$ radici \emph{distinte}
$\alpha_1, \dots, \alpha_k$. Allora esistono costanti
$c_1, \dots, c_k$ tali che per ogni $n \geq 0$
\[
  a_n = c_1\, \alpha_1^{\,n} + c_2\, \alpha_2^{\,n} + \dots + c_k\, \alpha_k^{\,n} .
\]
Le costanti $c_i$ sono univocamente determinate dalle condizioni iniziali.
\end{teorema}

(Senza dimostrazione.)

\paragraph{Procedura per trovare la forma chiusa.}
\begin{enumerate}
  \item Scrivere il polinomio caratteristico $p(t)$.
  \item Trovare le sue $k$ radici distinte $\alpha_1, \dots, \alpha_k$
        (con il teorema del fattore).
  \item Scrivere $a_n = c_1 \alpha_1^{\,n} + \dots + c_k \alpha_k^{\,n}$.
  \item Imporre le condizioni iniziali $a_0, \dots, a_{k-1}$: si ottiene un
        sistema lineare di $k$ equazioni nelle incognite $c_1, \dots, c_k$.
  \item Risolvere il sistema e sostituire i valori trovati.
\end{enumerate}

\begin{esempio}
Sia $a_n = 5\, a_{n-1} - 6\, a_{n-2}$ con $a_0 = 0$, $a_1 = 1$. Il polinomio
caratteristico è $p(t) = t^2 - 5t + 6 = (t - 2)(t - 3)$, con radici distinte
$2$ e $3$. Quindi
\[
  a_n = c_1\, 2^n + c_2\, 3^n .
\]
Le condizioni iniziali danno il sistema
\[
  \begin{cases}
    a_0 = c_1 + c_2 = 0 \\
    a_1 = 2c_1 + 3c_2 = 1
  \end{cases}
  \quad\Longrightarrow\quad c_2 = 1,\ c_1 = -1 ,
\]
perciò
\[
  a_n = 3^n - 2^n .
\]
Verifica: $a_2 = 5 \cdot 1 - 6 \cdot 0 = 5 = 3^2 - 2^2$ e
$a_3 = 5 \cdot 5 - 6 \cdot 1 = 19 = 3^3 - 2^3$.
\end{esempio}

\begin{esempio}[Fibonacci]
Per $F_n = F_{n-1} + F_{n-2}$, $F_0 = 0$, $F_1 = 1$, il polinomio
caratteristico è $t^2 - t - 1$, con radici distinte
\[
  \varphi = \frac{1 + \sqrt{5}}{2} , \qquad \psi = \frac{1 - \sqrt{5}}{2} .
\]
Allora $F_n = c_1\, \varphi^n + c_2\, \psi^n$ e
\[
  \begin{cases}
    c_1 + c_2 = 0 \\
    c_1\, \varphi + c_2\, \psi = 1
  \end{cases}
  \quad\Longrightarrow\quad
  c_1 (\varphi - \psi) = 1 ,\quad \varphi - \psi = \sqrt{5}
  \quad\Longrightarrow\quad c_1 = \frac{1}{\sqrt{5}},\ c_2 = -\frac{1}{\sqrt{5}} .
\]
Si ottiene la formula di Binet:
\[
  F_n = \frac{1}{\sqrt{5}}\left[\left(\frac{1 + \sqrt{5}}{2}\right)^{n}
        - \left(\frac{1 - \sqrt{5}}{2}\right)^{n}\right] .
\]
\end{esempio}

\begin{esempio}[grado 3]
Sia $a_n = 6 a_{n-1} - 11 a_{n-2} + 6 a_{n-3}$ con $a_0 = 3$, $a_1 = 6$,
$a_2 = 14$. Il polinomio caratteristico è
$p(t) = t^3 - 6t^2 + 11t - 6$. Provando i divisori di $6$ si trova
$p(1) = 0$, e con Ruffini $p(t) = (t - 1)(t^2 - 5t + 6) = (t-1)(t-2)(t-3)$.
Le radici $1, 2, 3$ sono distinte, quindi
\[
  a_n = c_1 + c_2\, 2^n + c_3\, 3^n ,
  \qquad
  \begin{cases}
    c_1 + c_2 + c_3 = 3 \\
    c_1 + 2c_2 + 3c_3 = 6 \\
    c_1 + 4c_2 + 9c_3 = 14
  \end{cases}
\]
Sottraendo la prima equazione dalla seconda e la seconda dalla terza si ha
$c_2 + 2c_3 = 3$ e $2c_2 + 6c_3 = 8$, da cui $c_3 = 1$, $c_2 = 1$, $c_1 = 1$.
Dunque
\[
  a_n = 1 + 2^n + 3^n .
\]
\end{esempio}

\subsubsection{Teorema B: radici con molteplicità}

\begin{teorema}[B]
Sia $(a_n)$ una successione ricorsiva lineare omogenea di grado $k$ il cui
polinomio caratteristico si scompone come
\[
  p(t) = (t - \alpha_1)^{m_1} (t - \alpha_2)^{m_2} \cdots (t - \alpha_r)^{m_r},
  \qquad m_1 + m_2 + \dots + m_r = k ,
\]
con $\alpha_1, \dots, \alpha_r$ distinte e $m_i$ la molteplicità di
$\alpha_i$. Allora
\[
  a_n = \sum_{i=1}^{r}
        \bigl(c_{i,0} + c_{i,1}\, n + \dots + c_{i,m_i - 1}\, n^{m_i - 1}\bigr)\,
        \alpha_i^{\,n} ,
\]
con $k$ costanti $c_{i,j}$ univocamente determinate dalle condizioni iniziali.
\end{teorema}

(Senza dimostrazione.)

\begin{osservazione}
Ogni radice $\alpha$ di molteplicità $m$ contribuisce con $m$ termini,
\[
  \alpha^n,\ n\,\alpha^n,\ n^2\alpha^n,\ \dots,\ n^{m-1}\alpha^n ,
\]
in totale $k$ costanti, quante le condizioni iniziali. Se tutte le molteplicità
sono $1$ si ritrova il Teorema A.
\end{osservazione}

\begin{esempio}[radice doppia]
Sia $a_n = 4\, a_{n-1} - 4\, a_{n-2}$ con $a_0 = 1$, $a_1 = 0$. Il polinomio
caratteristico è $p(t) = t^2 - 4t + 4 = (t - 2)^2$: un'unica radice $2$ di
molteplicità $2$. Quindi
\[
  a_n = (c_0 + c_1\, n)\, 2^n .
\]
Dalle condizioni iniziali:
\[
  a_0 = c_0 = 1 , \qquad a_1 = 2\,(c_0 + c_1) = 0 \Rightarrow c_1 = -1 ,
\]
perciò
\[
  a_n = (1 - n)\, 2^n .
\]
Verifica: $a_2 = 4 \cdot 0 - 4 \cdot 1 = -4 = (1-2)\cdot 2^2$ e
$a_3 = 4\cdot(-4) - 4 \cdot 0 = -16 = (1-3)\cdot 2^3$.
\end{esempio}

\begin{esempio}[radici di molteplicità diversa]
Sia $a_n = 5\, a_{n-1} - 8\, a_{n-2} + 4\, a_{n-3}$ con $a_0 = 2$, $a_1 = 5$,
$a_2 = 13$. Il polinomio caratteristico è $p(t) = t^3 - 5t^2 + 8t - 4$.
Poiché $p(1) = 0$, con Ruffini si ottiene
\[
  p(t) = (t - 1)(t^2 - 4t + 4) = (t - 1)(t - 2)^2 ,
\]
cioè la radice $1$ ha molteplicità $1$ e la radice $2$ ha molteplicità $2$.
Per il Teorema B
\[
  a_n = c + (d + e\, n)\, 2^n ,
  \qquad
  \begin{cases}
    c + d = 2 \\
    c + 2d + 2e = 5 \\
    c + 4d + 8e = 13
  \end{cases}
\]
Sottraendo la prima equazione dalla seconda e la seconda dalla terza:
$d + 2e = 3$ e $2d + 6e = 8$, cioè $d + 3e = 4$. Quindi $e = 1$, $d = 1$,
$c = 1$ e
\[
  a_n = 1 + (1 + n)\, 2^n .
\]
Verifica: $a_3 = 5 \cdot 13 - 8 \cdot 5 + 4 \cdot 2 = 33 = 1 + 4 \cdot 8$.
\end{esempio}

\end{document}
